\documentclass[11pt,a4paper]{article}
\usepackage{isabelle,isabellesym}
\usepackage{amsmath,amssymb}
\usepackage{hyperref}

\title{Vibesnipe: AGM Belief Revision meets Semi-Reliable Nashator}
\author{boxxy formalization}

\begin{document}

\maketitle

\begin{abstract}
We formalize the intersection of AGM belief revision theory and 
compositional game theory via Hedges-Capucci's semi-reliable nashator.
The \emph{vibesnipe} construction models multi-agent belief revision 
as a game where agents select revision operators from indeterministic 
options, reaching $\epsilon$-approximate Nash equilibria.

Key contributions:
\begin{itemize}
\item Isabelle/HOL formalization of AGM postulates (K*1--K*8)
\item Selection relations and the Nash product (nashator)
\item Semi-reliable variant for bounded rationality
\item Connection to Grove sphere systems and epistemic entrenchment
\item GF(3) trit conservation throughout
\end{itemize}
\end{abstract}

\tableofcontents

\section{Introduction}

The AGM framework \cite{agm1985} provides postulates for rational belief revision.
When epistemic entrenchment fails to be totally ordered, multiple revision 
operators become admissible---the Lindström-Rabinowicz indeterminism phenomenon.

Hedges and Capucci's compositional game theory \cite{hedges2021} provides 
selection functions and the Nash product for modeling multi-agent equilibria.
The \emph{semi-reliable} variant allows $\epsilon$-slack, modeling bounded 
rationality.

\textbf{Vibesnipe} = selecting from indeterministic revisions (vibes) with 
near-optimal precision (snipe).

\section{Theories}

\input{session}

\section{GF(3) Conservation}

All structures maintain GF(3) trit balance:
\begin{align}
\text{Selection } (+1) + \text{Nashator } (0) + \text{Verify } (-1) &\equiv 0 \pmod{3}
\end{align}

This ensures compositional coherence across the Cat\# equipment structure.

\bibliographystyle{plain}
\begin{thebibliography}{9}
\bibitem{agm1985} Alchourrón, Gärdenfors, Makinson. 
\emph{On the logic of theory change}. J. Symbolic Logic, 1985.

\bibitem{hedges2021} Capucci, Gavranović, Hedges, Rischel. 
\emph{Towards Foundations of Categorical Cybernetics}. ACT 2021.

\bibitem{capucci2022} Capucci. 
\emph{Diegetic Representation of Feedback in Open Games}. ACT 2022.

\bibitem{lindstrom1999} Lindström, Rabinowicz.
\emph{DDL unlimited: Dynamic doxastic logic for introspective agents}. 1999.
\end{thebibliography}

\end{document}
